\chapter{尺度、吸收与保留信息的导出}
\label{chap:scaling}

\section{为什么等价反射仍然强弱不同}

理想结构因子只描述晶体内部的相干散射。同一反射在不同扫描位置的观测，
还受入射通量、照射体积、偏振、Lorentz 因子和吸收影响。
若积分结果已经是完整轮廓强度，一个单独反射的前向模型可写成
\begin{equation}
 \operatorname{E}[y_j]=g_jI_{h(j)},\qquad
 g_j=k_j\,L_jP_j\,T_j>0.
 \label{eq:scaling-forward}
\end{equation}
这里 $L_jP_j$ 表示已声明约定的 Lorentz--偏振乘子，并非上一章的体素轮廓矩阵。
实现中的字段是 \code{lp_forward}；方向、单位和是否已经校正必须随数据保存。
完整轮廓积分之后再次除以可见比例，会把同一效应校正两次。

吸收来自光子沿路径的损失。线性吸收系数 $\mu$ 以 mm$^{-1}$ 表示时，
路径长度必须使用 mm。单一路径及有限晶体的平均透过率分别为
\begin{equation}
 T(\bm r)=\exp[-\mu\{L_{\mathrm{in}}(\bm r)+L_{\mathrm{out}}(\bm r)\}],
 \qquad
 \overline T=\frac{1}{V}\int_V T(\bm r)\,d^3\bm r.
 \label{eq:scaling-absorption}
\end{equation}
已知形状与面指数时，这给出物理吸收模型；缺少可靠形状时，多扫描等价反射可以约束经验方向响应。
两者解决的问题相近，但经验系数不应解释成直接测得的晶体形状或绝对吸收系数\cite{scaling2020}。

\section{现有方向模型到底包含什么}

沿用的 Ariadne 全局缩放采用正值指数模型：
\begin{equation}
 g_j=g_j^{\mathrm{known}}\exp(\bm x_j^T\bm\beta).
 \label{eq:scaling-exponential}
\end{equation}
当前代码支持每个扫描的零至二次角度特征，以及关闭吸收项或保留偶二阶方向项。
对共同晶体坐标系中的单位方向 $\bm v=(x,y,z)$，所用基为
\begin{equation}
 H_2(\bm v)=
 (xy,xz,yz,x^2-y^2,3z^2-1)^T,\qquad
 \bm a_j=H_2(\widehat{\bm s}_{\mathrm{in}})
          +H_2(\widehat{\bm s}_{\mathrm{out}}).
 \label{eq:scaling-harmonics}
\end{equation}
这是实二阶球谐空间的一组基，未采用正交归一化常数；系数会吸收该常数。
正值指数保证响应可用作尺度，但经过中心化的相对响应可以大于 1。
它不是必须小于 1 的绝对透过率，也不是任意阶球谐吸收模型。

\section{复合观测应在前向模型中缩放}

重叠组可以给出一个总强度、多个可分辨分量，或一组正交信息模态。
用 $D_b$ 表示这些观测对原贡献的线性映射，$E_b$ 将贡献连接到各自等价反射，则
\begin{equation}
 \bm\mu_b
 =D_b\,\operatorname{diag}\!\bigl(\bm g_b(\bm\beta)\bigr)E_b\bm I.
 \label{eq:scaling-composite}
\end{equation}
单一总量对应一行全 1 的 $D_b$；已辨识分量对应单位阵；压缩 Gaussian 观测则可直接使用上一阶段的 $R$。
这样可以使用不同扫描上的等价信息，而无须先给不可分辨像素强行拆分强度。
每个输入组保留自己的协方差，再统一白化；同一组内的相关分量不能当成互相独立的行。

各孪晶域共享尺度约定也很关键。逐域把平均强度调成一样，会吸收本来属于域体积分数的信息。
保留代码不进行这种独立域归一化；0.8 的结构桥接进一步用结构因子和域分数预测贡献，再进入原观测算子。

\section{尺度规范与变量投影}

强度和尺度存在可交换的乘数：$I\mapsto e^{-c}I$、$g\mapsto e^cg$
不改变预测。现有实现将贡献特征在全局中心化，固定相对尺度规范：
\begin{equation}
 \sum_j\bm x_j=0
 \quad\Longrightarrow\quad
 \sum_j\log\frac{g_j}{g_j^{\mathrm{known}}}=0.
 \label{eq:scaling-gauge}
\end{equation}
这只消除整体规范。若某个方向特征在每个等价反射组内都不变化，它仍可被该组未知强度吸收。
代码另检查组内中心化特征的秩，保留剩余零空间；缺少角度覆盖不会因添加正则项而自动变成实验信息。

给定非线性尺度 $\bm\beta$，有符号强度仍是线性未知量。
把所有白化观测拼成 $\bm d$，相应设计记为 $A(\bm\beta)$，可先求
\begin{equation}
 \widehat{\bm I}(\bm\beta)=A(\bm\beta)^+\bm d,
 \qquad
 \bm r(\bm\beta)=A\widehat{\bm I}-\bm d.
 \label{eq:scaling-varpro}
\end{equation}
这是变量投影\cite{golub1973}。在线性设计保持满列秩的邻域，完整残差导数为
\begin{equation}
 d\bm r=(1-AA^+)(dA)\widehat{\bm I}
          -A^{+T}(dA)^T\bm r.
 \label{eq:scaling-varpro-derivative}
\end{equation}
第二项在残差非零时不能省略。现有实现用 QR/SVD 和固定支撑分组减少重复计算，
避免解正规方程造成条件数平方；满秩条件失效时返回原交替路径，并继续报告秩与收敛证据。
这里的有符号线性解仍然保留负强度估计，不为获得好看的合并统计而截零。

\section{误差传播与导出不能脱节}

设校正后的可测总量为 $S=\bm f^T\bm I$。
如果尺度也从这批强度估计出来，强度与校正因子通常相关，因此
\begin{equation}
 \operatorname{Var}(S)=
 \bm f^TC_I\bm f+\bm I^TC_f\bm I
 +2\bm f^TC_{If}\bm I.
 \label{eq:scaling-cross-covariance}
\end{equation}
最后的同数据交叉项不能因为实现方便而省略。
沿用缩放代码保存联合参数协方差和观测灵敏度，校正阶段据此传播误差。
若像素只给出 $I_1+I_2$，而两分量需要不同的校正因子，仅凭该总量不能得到校正后的和；
应保留复合前向模型，不能用一次结构预测拆分冒充两个新的实验测量。

HKLF 5 是交换格式，不足以装下完整的分量与跨观测协方差。
现有导出器只接受有完整贡献者记录、同相证据和真实校正来源的观测组。
为适应固定小数字段，它可以对所有行共同变换
\begin{equation}
 I'_j=aI_j,\qquad \sigma'_j=a\sigma_j,\qquad a>0,
 \label{eq:scaling-export}
\end{equation}
从而在小数舍入前保持 $I/\sigma$。负强度照常保留；若字段动态范围无法同时容纳有效误差条，
程序保留主数据并报错，不把过小的 sigma 换成任意常数。
非负计数的剖面区间也不能通过“区间宽度除以 2”伪装成 Gaussian sigma。
主 JSON/JSONL 中的似然、协方差和行映射应与 HKL 同时交付。

\begin{table}[!htbp]
 \centering\small
 \caption{沿用的缩放和导出实现；路径前缀为
 \code{skills/ariadne-reduction/scripts/}。}
 \begin{tabularx}{\textwidth}{@{}p{0.40\textwidth}X@{}}
  \toprule
  模块 & 数学对象及作用 \\
  \midrule
  \code{global_scale.py} & \code{harmonics} 与 \code{prepare} 构造共同方向特征、等价关系和观测块。\\
  \code{composite_scaling.py} & \code{fit_global} 联合估计尺度与强度，检查规范、秩与曲率。\\
  \code{variable_projection.py} & 线性强度消元及完整的残差导数。\\
  \code{correct_overlap_groups.py} & 包含同数据交叉项的校正误差传播。\\
  \code{hklf5_stream.py} & 两遍流式格式转换，记录舍入误差与信息损失。\\
  \bottomrule
 \end{tabularx}
\end{table}

本章解释保留的 0.7 数值能力。0.8 新增的是原生像素似然和结构前向桥接，
没有把经验二阶吸收模型重新命名为完整形状吸收，也没有以合成测试代替真实原始帧的缩放验证。
